\section{Generation-2 FAST surrogate factorization}
\label{sec:generation2-fast}

The Generation-2 FAST surrogate is the concrete trainable realization of the
physics-factored architecture in Section~\ref{sec:physics-factored}.  Its
purpose is to preserve the physical invariants already enforced by the vNext
REFERENCE stack while removing three avoidable sources of learned-model bias:
a one-pass many-body graph, an unnecessarily positive residual resistance, and
a spatial target that depends on upstream Port prediction error.

\subsection{Shared immutable scene partition}

Teacher scene index $j$ is mapped by a fixed hash $h(j,s)$, where $s$ is one
split seed, into exactly one of
\[
  \mathcal D_{\rm train},\qquad
  \mathcal D_{\rm val},\qquad
  \mathcal D_{\rm test}.
\]
Membership is independent of the current dataset length.  Port and spatial
training use the same partition.  Active-learning additions may enter only
$\mathcal D_{\rm train}$; validation and test membership remain immutable.

\subsection{Repeated physics-initialized scene interactions}

Let $c_i^{(0)}$ and $p_a^{(0)}$ be normalized coil and package node embeddings.
Relative-pose edges also contain cheap dimensionless physical descriptors such
as inverse-distance factors, normalized package volume, electric/loss contrast,
and magnetic contrast.  Generation-2 performs a fixed number $K$ of shared
interaction rounds,
\[
 m_{cc,i}^{(r)}
 =\mathop{\mathrm{Agg}}_{j\ne i}
 M_{cc}(c_i^{(r)},c_j^{(r)},e_{ij}),
\]
\[
 m_{pp,a}^{(r)}
 =\mathop{\mathrm{Agg}}_{b\ne a}
 M_{pp}(p_a^{(r)},p_b^{(r)},e_{ab}),
\]
\[
 m_{cp,a}^{(r)}
 =\mathop{\mathrm{Agg}}_i
 M_{cp}(p_a^{(r)},c_i^{(r)},e_{ia}),
\]
\[
 m_{pc,i}^{(r)}
 =\mathop{\mathrm{Agg}}_a
 M_{pc}(c_i^{(r)},p_a^{(r)},e_{ia}).
\]
Residual node updates are
\[
 c_i^{(r+1)}
 =\operatorname{Norm}\!\left(
 c_i^{(r)}+U_c(c_i^{(r)},m_{cc,i}^{(r)},m_{pc,i}^{(r)})
 \right),
\]
\[
 p_a^{(r+1)}
 =\operatorname{Norm}\!\left(
 p_a^{(r)}+U_p(p_a^{(r)},m_{pp,a}^{(r)},m_{cp,a}^{(r)})
 \right).
\]
The first implementation uses $K=3$.  Increasing MLP depth inside a single edge
map is not considered a substitute for increasing this interaction order.

\subsection{Baseline-conditioned passive resistance}

Let $R_0\PSD$ be the analytic conductor/background baseline resistance and let
$B=R_0^{1/2}$ be a numerically stabilized PSD square root.  The neural decoder
predicts a real symmetric dimensionless log-correction $S=S^\T$.  Define
\[
 M=\exp(S)\SPD,
 \qquad
 R=BMB^\T\PSD.
\]
This parameterization has three required properties:
\begin{enumerate}
\item $S=0$ reproduces the analytic baseline exactly;
\item the final resistance is passive for every network output;
\item $R-R_0$ is not constrained to be positive semidefinite.
\end{enumerate}
The third property is essential because passivity constrains the final total
operator, not the sign of the package-induced correction to a package-free
baseline.

The reactive decoder remains signed and reciprocal,
\[
 X=X_0+\Delta X,
 \qquad
 \Delta X=\Delta X^\T.
\]

\subsection{Power-closing dissipative channels}

Let $\widetilde D_k\PSD$ be raw learned channel operators.  Their raw sum is
\[
 \widetilde R=\sum_k\widetilde D_k.
\]
A low-dimensional congruence $T$ is computed such that
\[
 T\widetilde R T\adj=R.
\]
The returned channels are
\[
 D_k=T\widetilde D_kT\adj\PSD,
 \qquad
 \sum_kD_k=R.
\]
Thus passivity and exact Port power closure are independent of training error.
The initial Generation-2 migration retains the existing teacher channel family
(one conductor channel per coil plus one aggregate electric-environment
channel) so the energy-v3 teacher cache remains reusable.

\subsection{Teacher-channel canonical spatial shape}

The spatial surrogate learns geometry-dependent loss shape rather than upstream
Port amplitude error.  For teacher channel $D_k^\star$ and teacher local loss
operator $H_k^\star(\xi)$, define a normalized shape representation by a
congruence that enforces
\[
 \int H_{k,\theta}^{\rm canon}(\xi)\,\dd\mu
 =D_k^\star.
\]
Spatial training compares this teacher-channel-normalized prediction directly
with $H_k^\star(\xi)$.  The teacher target is therefore independent of the Port
surrogate prediction.

At inference the same raw spatial shape is normalized instead to the predicted
Port channel $D_{k,\theta}$,
\[
 \int H_{k,\theta}^{\rm run}(\xi)\,\dd\mu
 =D_{k,\theta}.
\]
This preserves deployed Port--Spatial power closure while separating Port
channel error from spatial shape error during training and validation.

\subsection{Spatial-specific scene refinement}

The spatial network receives the frozen Port latent but does not treat it as a
sufficient statistic.  It also receives normalized intrinsic and relative scene
features and performs a small spatial-specific interaction refinement before
query-point decoding.  This permits local interface, anisotropy, and boundary
information to remain available even when it is not required by the Port head.

\subsection{Loss decomposition and model selection}

Port optimization reports separate symmetric relative losses
\[
 \mathcal L_R,\qquad
 \mathcal L_X,\qquad
 \mathcal L_D,
\]
and minimizes
\[
 \mathcal L_{\rm port}
 =w_R\mathcal L_R+w_X\mathcal L_X+w_D\mathcal L_D.
\]
Because the spatial stage depends directly on dissipative channels, the initial
Generation-2 configuration uses $w_D>w_R=w_X$ while still reporting every
component independently.

Spatial training and selection use the same canonical shape metric.  Full
Port--Spatial end-to-end error is reported separately and cannot be mistaken
for a same-objective train/validation pair.

\subsection{Cache and artifact boundary}

The physical energy-v3 teacher cache is representation-independent and remains
valid.  Generation-2 re-encodes the stored scene at training time.  Neural
checkpoints and artifacts are not representation-independent: they carry a new
generation/schema identifier and cannot resume Generation-1 optimizer state.

This separation allows a clean architecture experiment on the same physical
teacher scenes before changing the sampling distribution.